% Required packages (add to your preamble if not already present):
%   \usepackage[skins,breakable]{tcolorbox}
%   \usepackage{enumitem}
%   \usepackage{booktabs}
%   \usepackage{xcolor}

\noindent This appendix summarizes a programmatic encoding of machine-checkable provisions from the shared entry rulebook, as used by \texttt{detect.py} to check \texttt{analysis.md} and kernel files. Check types include \texttt{section\_exists}, \texttt{pattern}, \texttt{field\_value}, and \texttt{cross\_check}.\par\smallskip
\begin{tcolorbox}[breakable,enhanced,colback=white,colframe=gray!50,boxrule=0.8pt,arc=4pt,left=8pt,right=8pt,top=6pt,bottom=6pt,title={\normalsize\bfseries NeuroDoc Semantic Rulebook},coltitle=black,colbacktitle=gray!15,titlerule=0.5pt]
\noindent\rule{\linewidth}{0.4pt}\par\smallskip
\noindent{\bfseries\small Hard Constraints (HC)}
{\footnotesize\ ---\ enforced in generation pipeline; not detectable by \texttt{detect.py}}\par\smallskip
\noindent{\footnotesize\textbf{HC-1.}~Data fields in File Scan Result (Event ID, Count, Info table values, Kernel SHA256 fingerprint) must come from kernel scan program output; LLM must not invent or alter these fields. Semantic fields (Event Type, Meaning) are generated by LLM in an isolated inference step but must be based on actual evidence from dataset files; fabrication is not allowed.}\par\smallskip
\noindent{\footnotesize\textbf{HC-2.}~Event Meaning and Type inference must be based on actual evidence in the dataset (events.tsv, sidecar JSON, README, etc.). LLM may consult any context information, but must not use unrelated reference papers, domain conventions, or BCI competition code systems as inference basis (e.g., must not infer four motor-imagery classes solely because event IDs are 1-4). Once an inference result is fixed it becomes a locked value; downstream pipeline must not modify it.}\par\smallskip
\noindent{\footnotesize\textbf{HC-3.}~Meta Area fields (Channels, Sampling Rate) should be cross-referenced against the locator cache CSV. LLM may fill narrative gaps but must not override program-derived values.}\par\smallskip
\noindent{\footnotesize\textbf{HC-4.}~LLM or automation may only set Status to: wait\_manual\_review -- default output when no explicit structural kernel gap is exposed wait\_kernel\_update -- explicit structural kernel gap or scan failure "completed" and "hallucination" are reserved for human reviewers only.}\par\smallskip
\noindent{\footnotesize\textbf{HC-5.}~Kernel must NOT remap original event IDs. Prefix-based enrichment is allowed (e.g., subject\_attr-event\_id, session\_attr-event\_id). Original event IDs must pass through unchanged.}\par\smallskip
\noindent{\footnotesize\textbf{HC-6.}~Kernel channel remapping is only allowed when explicit evidence exists in the dataset directory (companion file, README, sidecar) mapping ambiguous/numeric channels to named channels. Add a comment in kernel citing the evidence file.}\par\smallskip
\noindent{\footnotesize\textbf{HC-7.}~Full scan count should be minimized to avoid re-scanning the same dataset. The pipeline should run a consistency check before final output to ensure document content matches scan results.}\par\smallskip
\noindent\rule{\linewidth}{0.4pt}\par\smallskip
\noindent{\bfseries\small Semantic Check Rules}\par\smallskip
\smallskip\noindent{\footnotesize\bfseries Group~A:~Document Structure}\par\smallskip
\begingroup\footnotesize
\begin{tabular}{@{}p{0.22\linewidth}ll p{0.46\linewidth}@{}}
\toprule
\textbf{Rule ID} & \textbf{Sev.} & \textbf{Target} & \textbf{Description} \\
\midrule
\texttt{has\_\allowbreak{}paradigm} & \textbf{[E]} & \texttt{analysis\_md} & analysis.md must contain \#\# Paradigm section. \\
\addlinespace[1pt]
\texttt{has\_\allowbreak{}file\_\allowbreak{}scan\_\allowbreak{}result} & \textbf{[E]} & \texttt{analysis\_md} & analysis.md must contain \#\# File Scan Result section. \\
\addlinespace[1pt]
\texttt{has\_\allowbreak{}tasks} & \textbf{[E]} & \texttt{analysis\_md} & analysis.md must contain \#\# Tasks section. \\
\addlinespace[1pt]
\texttt{has\_\allowbreak{}data\_\allowbreak{}quality\_\allowbreak{}notes} & \textbf{[E]} & \texttt{analysis\_md} & analysis.md must contain \#\# Data Quality Notes section. \\
\addlinespace[1pt]
\texttt{has\_\allowbreak{}kernel\_\allowbreak{}responsibilities} & \textbf{[E]} & \texttt{analysis\_md} & analysis.md must contain \#\# Kernel Responsibilities section. \\
\addlinespace[1pt]
\texttt{has\_\allowbreak{}event\_\allowbreak{}structure} & \textbf{[E]} & \texttt{analysis\_md} & \#\# File Scan Result must contain \#\#\# Event Structure subsection. \\
\addlinespace[1pt]
\texttt{has\_\allowbreak{}info\_\allowbreak{}subsection} & \textbf{[E]} & \texttt{analysis\_md} & \#\# File Scan Result must contain \#\#\# Info subsection (rulebook Article 57.1). \\
\addlinespace[1pt]
\texttt{no\_\allowbreak{}stim\_\allowbreak{}label\_\allowbreak{}summary} & \textbf{[E]} & \texttt{analysis\_md} & STIM Label Summary is a deprecated legacy section and must not appear. \\
\addlinespace[1pt]
\texttt{section\_\allowbreak{}order\_\allowbreak{}correct} & \textbf{[E]} & \texttt{analysis\_md} & Top-level sections must follow order: Paradigm -> File Scan Result -> Tasks -> Data Quality Notes -> Kernel Responsibilities. \\
\addlinespace[1pt]
\texttt{no\_\allowbreak{}extra\_\allowbreak{}top\_\allowbreak{}level\_\allowbreak{}sections} & \textbf{[E]} & \texttt{analysis\_md} & Only the 5 required \#\# sections are allowed. No extra top-level headings (rulebook Article 30.2). \\
\addlinespace[1pt]
\bottomrule
\end{tabular}
\endgroup
\smallskip\noindent{\footnotesize\bfseries Group~B:~Field Legal Values}\par\smallskip
\begingroup\footnotesize
\begin{tabular}{@{}p{0.22\linewidth}ll p{0.46\linewidth}@{}}
\toprule
\textbf{Rule ID} & \textbf{Sev.} & \textbf{Target} & \textbf{Description} \\
\midrule
\texttt{status\_\allowbreak{}valid\_\allowbreak{}value} & \textbf{[E]} & \texttt{analysis\_md} & Status must be one of: wait\_kernel\_update, wait\_manual\_review, completed, hallucination (rulebook Article 42.2). \\
\addlinespace[1pt]
\texttt{super\_\allowbreak{}meta\_\allowbreak{}required\_\allowbreak{}valid\_\allowbreak{}value} & \textbf{[E]} & \texttt{analysis\_md} & Super Meta Required must be exactly `Yes` or `No` (rulebook Article 37). \\
\addlinespace[1pt]
\texttt{category\_\allowbreak{}valid\_\allowbreak{}value} & \textbf{[E]} & \texttt{analysis\_md} & Category must be one of 6 valid values (rulebook Article 25): Complex Applied / High-level Mental / Action/Output / Sensory/Response / Biological State / Others. \\
\addlinespace[1pt]
\texttt{label\_\allowbreak{}route\_\allowbreak{}no\_\allowbreak{}deprecated} & \textbf{[E]} & \texttt{analysis\_md} & Label Route must use current values event, misc, or info. Historical values Annotation/STIM/MISC are not allowed (rulebook Article 77). \\
\addlinespace[1pt]
\texttt{event\_\allowbreak{}type\_\allowbreak{}no\_\allowbreak{}deprecated} & \textbf{[E]} & \texttt{analysis\_md} & Event Structure Type must use event or admin only. Historical short/long/rest values are not allowed (rulebook Article 46). \\
\addlinespace[1pt]
\texttt{kernel\_\allowbreak{}fingerprint\_\allowbreak{}exists} & \textbf{[E]} & \texttt{analysis\_md} & \#\# File Scan Result must begin with a kernel fingerprint blockquote containing Kernel SHA256 (rulebook Article 42.5). \\
\addlinespace[1pt]
\texttt{no\_\allowbreak{}placeholder\_\allowbreak{}doi} & \textbf{[W]} & \texttt{analysis\_md} & Evidence DOIs must not be placeholder strings (TODO / N/A / none / unknown / placeholder). \\
\addlinespace[1pt]
\bottomrule
\end{tabular}
\endgroup
\smallskip\noindent{\footnotesize\bfseries Group~C:~Semantic Consistency}\par\smallskip
\begingroup\footnotesize
\begin{tabular}{@{}p{0.22\linewidth}ll p{0.46\linewidth}@{}}
\toprule
\textbf{Rule ID} & \textbf{Sev.} & \textbf{Target} & \textbf{Description} \\
\midrule
\texttt{super\_\allowbreak{}meta\_\allowbreak{}matches\_\allowbreak{}tasks} & \textbf{[E]} & \texttt{analysis\_md} & Super Meta Required must match actual Meta Level in Tasks: Yes if any task has Meta Level=super, No if all raw (rulebook Article 37-38). \\
\addlinespace[1pt]
\texttt{target\_\allowbreak{}events\_\allowbreak{}in\_\allowbreak{}event\_\allowbreak{}structure} & \textbf{[E]} & \texttt{analysis\_md} & All event IDs referenced in Tasks Target Definition tables must appear in \#\#\# Event Structure. \\
\addlinespace[1pt]
\texttt{classical\_\allowbreak{}task\_\allowbreak{}has\_\allowbreak{}evidence} & \textbf{[E]} & \texttt{analysis\_md} & Every Classical Task must have a non-empty Evidence field (rulebook Articles 62 and 79). \\
\addlinespace[1pt]
\texttt{info\_\allowbreak{}completeness\_\allowbreak{}valid} & \textbf{[E]} & \texttt{analysis\_md} & \#\#\# Info table Completeness column must use: complete, partial (x/n), incomplete, or missing (rulebook Article 57.3). \\
\addlinespace[1pt]
\texttt{kernel\_\allowbreak{}sha256\_\allowbreak{}matches\_\allowbreak{}file} & \textbf{[E]} & \texttt{analysis\_md} & If Kernel SHA256 is a concrete hash, it must match the effective kernel file used for scanning. \texttt{N/A} is allowed when scanning does not require a dataset-specific kernel file (for example generic/default parser cases). \\
\addlinespace[1pt]
\texttt{no\_\allowbreak{}auto\_\allowbreak{}removed\_\allowbreak{}tasks} & \textbf{[W]} & \texttt{analysis\_md} & Tasks section must not contain [AUTO-REMOVED] blocks. These indicate guard-removed hallucinations; dataset needs regeneration. \\
\addlinespace[1pt]
\texttt{no\_\allowbreak{}mi\_\allowbreak{}label\_\allowbreak{}on\_\allowbreak{}phase\_\allowbreak{}marker} & \textbf{[E]} & \texttt{analysis\_md} & Event Structure must not assign MI class labels to events whose Meaning is a trial phase marker (rest/ready/task/interval). See HC-2. \\
\addlinespace[1pt]
\texttt{no\_\allowbreak{}hallucinated\_\allowbreak{}bcic\_\allowbreak{}classes} & \textbf{[W]} & \texttt{analysis\_md} & Must not assign BCIC IV 2a class codes (769/770/771/772) as event labels without explicit sidecar evidence. \\
\addlinespace[1pt]
\texttt{kr\_\allowbreak{}has\_\allowbreak{}super\_\allowbreak{}meta} & \textbf{[E]} & \texttt{analysis\_md} & \#\# Kernel Responsibilities must contain > \textbf{Super Meta Required}: blockquote field (rulebook Articles 36-38). \\
\addlinespace[1pt]
\texttt{kr\_\allowbreak{}has\_\allowbreak{}status} & \textbf{[E]} & \texttt{analysis\_md} & \#\# Kernel Responsibilities must contain > \textbf{Status}: blockquote field (rulebook Articles 36 and 42.1). \\
\addlinespace[1pt]
\texttt{kr\_\allowbreak{}only\_\allowbreak{}blockquote\_\allowbreak{}fields} & \textbf{[W]} & \texttt{analysis\_md} & \#\# Kernel Responsibilities must contain only the two blockquote fields (Super Meta Required, Status), no extra prose (rulebook Articles 36 and 42.4). \\
\addlinespace[1pt]
\texttt{dqn\_\allowbreak{}bullet\_\allowbreak{}format} & \textbf{[W]} & \texttt{analysis\_md} & \#\# Data Quality Notes content must use `- ` bullet format (rulebook Article 83). \\
\addlinespace[1pt]
\texttt{task\_\allowbreak{}type\_\allowbreak{}no\_\allowbreak{}normal} & \textbf{[E]} & \texttt{analysis\_md} & Task headings must be `Classical Task` or `Potential Task` only. `Normal Task` is a legacy label and is not allowed. \\
\addlinespace[1pt]
\texttt{task\_\allowbreak{}order\_\allowbreak{}classical\_\allowbreak{}first} & \textbf{[E]} & \texttt{analysis\_md} & All Classical Tasks must appear before all Potential Tasks in \#\# Tasks section. \\
\addlinespace[1pt]
\texttt{task\_\allowbreak{}numbering\_\allowbreak{}from\_\allowbreak{}one} & \textbf{[E]} & \texttt{analysis\_md} & Task numbering must start from 1 within each type and increase consecutively. \\
\addlinespace[1pt]
\texttt{task\_\allowbreak{}dataset\_\allowbreak{}dispatchable} & \textbf{[E]} & \texttt{analysis\_md} & Every Task block must be accepted by dataset\_api.build\_dataset\_from\_task(dry\_run=True). \\
\addlinespace[1pt]
\texttt{target\_\allowbreak{}def\_\allowbreak{}schema\_\allowbreak{}matches\_\allowbreak{}route} & \textbf{[E]} & \texttt{analysis\_md} & Target Definition table header must match the Label Route x Task Type matrix (rulebook Article 77.4). \\
\addlinespace[1pt]
\texttt{single\_\allowbreak{}target\_\allowbreak{}per\_\allowbreak{}task} & \textbf{[E]} & \texttt{analysis\_md} & For misc/info tasks, the first Target Definition column must contain exactly one unique source value. This rule does not restrict event tasks with multiple Event ID rows (rulebook Articles 77.3 and 77.5). \\
\addlinespace[1pt]
\texttt{classification\_\allowbreak{}task\_\allowbreak{}has\_\allowbreak{}multiple\_\allowbreak{}classes} & \textbf{[E]} & \texttt{analysis\_md} & Classification tasks must have >= 2 distinct y values in Target Definition (rulebook Article 77.7). \\
\addlinespace[1pt]
\texttt{classification\_\allowbreak{}y\_\allowbreak{}zero\_\allowbreak{}indexed\_\allowbreak{}sequential} & \textbf{[E]} & \texttt{analysis\_md} & Event-route classification tasks must use 0-based sequential integer y values: the distinct y set must equal set(range(N)). Many-to-one mapping is allowed (e.g. A->0, B->1, C->1 is valid; \{0,1,3\} or \{1,2\} are not). Misc/info routes use Y Mapping Notes instead of a y column. \\
\addlinespace[1pt]
\texttt{derived\_\allowbreak{}misc\_\allowbreak{}requires\_\allowbreak{}super\_\allowbreak{}meta} & \textbf{[E]} & \texttt{analysis\_md} & Tasks with Label Route = misc whose Target Definition source is not listed as native in MISC Label Summary must have Meta Level = super (rulebook Article 77.8). \\
\addlinespace[1pt]
\texttt{segmentation\_\allowbreak{}requires\_\allowbreak{}empty\_\allowbreak{}target\_\allowbreak{}events} & \textbf{[E]} & \texttt{analysis\_md} & Epoch Method = epoch\_by\_segmentation\_hdf5 requires target\_events = [] (rulebook Articles 45 and 72). \\
\addlinespace[1pt]
\texttt{epoch\_\allowbreak{}parameters\_\allowbreak{}complete} & \textbf{[E]} & \texttt{analysis\_md} & Epoch Parameters must list the minimum required keys for the declared Epoch Method (rulebook Articles 69-72). \\
\addlinespace[1pt]
\texttt{manual\_\allowbreak{}review\_\allowbreak{}blank} & \textbf{[W]} & \texttt{analysis\_md} & Completed documents must record at least one reviewer name. \texttt{wait\_manual\_review} documents may be blank or may retain reviewer names from later manual rework. \\
\addlinespace[1pt]
\bottomrule
\end{tabular}
\endgroup
\smallskip\noindent{\footnotesize\bfseries Group~D:~Kernel Code Constraints}\par\smallskip
\begingroup\footnotesize
\begin{tabular}{@{}p{0.22\linewidth}ll p{0.46\linewidth}@{}}
\toprule
\textbf{Rule ID} & \textbf{Sev.} & \textbf{Target} & \textbf{Description} \\
\midrule
\texttt{kernel\_\allowbreak{}has\_\allowbreak{}kernel\_\allowbreak{}id} & \textbf{[E]} & \texttt{kernel} & Every kernel file must define KERNEL\_ID. \\
\addlinespace[1pt]
\texttt{kernel\_\allowbreak{}has\_\allowbreak{}apply} & \textbf{[E]} & \texttt{kernel} & Every kernel file must implement apply(). \\
\addlinespace[1pt]
\texttt{kernel\_\allowbreak{}no\_\allowbreak{}shell\_\allowbreak{}commands} & \textbf{[E]} & \texttt{kernel} & Kernel must not use os.system, subprocess.call, or shell=True. \\
\addlinespace[1pt]
\texttt{kernel\_\allowbreak{}no\_\allowbreak{}file\_\allowbreak{}writes} & \textbf{[E]} & \texttt{kernel} & Kernel must not write or delete files. \\
\addlinespace[1pt]
\texttt{kernel\_\allowbreak{}no\_\allowbreak{}event\_\allowbreak{}remapping} & \textbf{[W]} & \texttt{kernel} & Kernel must not hard-remap original event IDs. Prefix enrichment (attr-event\_id) is allowed; replacing/renaming event IDs is not (HC-5). \\
\addlinespace[1pt]
\bottomrule
\end{tabular}
\endgroup
\end{tcolorbox}